\section{Introduction}
\label{sec:intro}

Genomic language models are distributed pretrained and fine-tuned locally. DNABERT
\cite{ji2021dnabert}, DNABERT-2 \cite{zhou2024dnabert2}, HyenaDNA \cite{nguyen2023hyenadna} and the
Nucleotide Transformer \cite{dalla2024nucleotide} are released as checkpoints that a group adapts to
its own sequences, and increasingly those sequences are cohort data. The adapted checkpoint is then
shared or published, and whatever it retained of the cohort travels with it.

What it retains matters more here than in text. A variant carried by a handful of people is close to
an identifier, and a model that completes the flanking sequence with the carrier's allele has
disclosed a genotype at a known locus rather than reproduced a string --- the adversary is not
guessing a sequence but resolving an allele whose position it already knows
\cite{li2013membership,tramer2015bounded}.

\cite{nemecek2026quantifying} established that the problem is real: four models --- a 12.9M
character-level transformer, DNABERT-2, HyenaDNA and Evo at 7B --- on four datasets, with large
differences in how much of the fine-tuning data comes back, attributed to ``architecture''. That is
where the question stops, because those four differ in architecture, parameter count by a factor of
\num{500}, tokenisation, training objective, pretraining corpus \emph{and} fine-tuning strategy all
at once --- ``architecture'' is a bundle, and a practitioner choosing a model cannot act on it. The closest result in text,
\cite{kharitonov2021bpe}, showed that a larger subword vocabulary memorises more and conjectured
that the shorter token sequences are the reason, but in text the two cannot be separated.

DNA separates them, for three reasons that belong to the alphabet rather than to our setup. It is
closed and four-valued, so a tokeniser is a pure re-grouping of the same information; uniform
synthetic DNA has an exactly known chance level of \SI{2.000}{\bit\per\nucleotide}, which calibrates
every measurement against a known zero; and token length can be varied with the token-sequence
length held fixed. We therefore hold the backbone, the data, the optimiser and the schedule fixed,
vary how DNA is split into tokens and whether the model predicts the next piece or fills in masked
pieces, and then check the resulting claims on the released checkpoints. Our contributions:

\begin{enumerate}\itemsep1pt
  \item \textbf{A controlled decomposition.} One fixed \SI{12.9}{\mega\nothing}-parameter backbone;
        only the tokeniser, the objective and the position encoding change. Across 24 runs the
        objective accounts for \NUM{67\%} of the variance in memorisation, the tokeniser for
        \NUM{12\%}, their interaction for \NUM{21\%} and the seed for \NUM{0\%}.
  \item \textbf{The objective decides whether records come back at all.} On all five datasets every
        masked cell reproduces \NUM{0\%} of the planted records and every coarse-tokeniser
        next-token cell reproduces \NUM{100\%}, with the masked cells at the \emph{better} held-out
        loss --- under the prefix-conditioned protocol and under the edge-conditioned infilling
        protocol that recovers up to three times more from bidirectional models
        \cite{wang2026infilling}. The released DNABERT checkpoints behave the same way under their
        own masked objective.
  \item \textbf{Monotonicity in average token length, with a mechanism.} Within an objective,
        memorisation rises with nucleotides per token on all five datasets. It is not the number of
        token slots, which predicts the opposite spread, and not sequence length, which we hold
        fixed; it is how often those tokens occur elsewhere in the corpus.
  \item \textbf{Utility gives no warning.} A validation loss \NUM{0.004}--\NUM{0.111}~\bits{} from
        optimal separates \NUM{0\%} from \NUM{51}--\NUM{100\%} verbatim reproduction; four of six
        GUE tasks \cite{zhou2024dnabert2} move by at most \NUM{0.037} MCC across that transition and
        the other two move \emph{up}; selecting by downstream score picks the leakiest checkpoint for
        five of six models.
  \item \textbf{Real haplotypes, including records seen once.} With canaries that are real
        chromosome-22 segments carrying rare 1000~Genomes variants \cite{thousandgenomes2015}, every
        released model returns the carrier's allele for \NUM{92}--\NUM{100\%} of variants in a
        segment seen sixteen times against a \NUM{23}--\NUM{27\%} control; DNABERT-2 reaches
        \NUM{85\%} for a segment seen \emph{once} and reproduces \NUM{58\%} of them verbatim, where a
        membership-ranking audit of the same checkpoints reads as safe.
  \item \textbf{Two settings, in order of effect.} Prefer a masked objective wherever generative
        output is not needed; tokenise finely as a second line, which delays first extraction by
        about an order of magnitude in rounds but does not prevent it.
\end{enumerate}

We do not claim that any setting makes memorisation impossible --- the character-level tokeniser,
our tightest cell, reaches complete reproduction when trained long enough, and we report the round.
We compare at matched held-out loss and equal exposure, not at matched compute. We study one human
chromosome and models up to \SI{117}{\mega\nothing} parameters, with canaries we plant rather than
discover, and an adversary assumed to hold the flanking 48\nt{}.
